\documentclass[11pt]{article}
\usepackage[a4paper,margin=1in]{geometry}
\usepackage{amsmath,amssymb,amsthm}
\usepackage{graphicx}
\usepackage[hidelinks]{hyperref}
\setlength{\emergencystretch}{3em}

\newtheorem{theorem}{Theorem}
\newtheorem{lemma}{Lemma}
\newtheorem{remark}{Remark}

\newcommand{\D}{\mathsf D}
\newcommand{\Lam}{\mathsf{La}}
\newcommand{\WB}{\mathsf{WB}}
\newcommand{\symdiff}{\triangle}
\newcommand{\diam}{\operatorname{diam}}
\newcommand{\dist}{\operatorname{dist}}

\title{Diamonds Do Not Embed Uniformly into Lamplighters on Weighted Binary Trees}
\author{Automatic math-discovery packet}
\date{July 18, 2026}

\begin{document}
\maketitle

\section*{Source Question}

The source paper is Charlotte Melby and Beata Randrianantoanina,
\emph{Bi-Lipschitz embedding properties of lamplighter graphs on weighted and
unweighted trees}, arXiv:2509.25421. After proving that diamonds do not embed
uniformly into lamplighter graphs on the subdivided trees \(W_{\theta,n}\) for
\(\theta\geq 3\), the authors ask whether the same theorem is true for the
weighted complete binary trees:

\begin{quote}
We do not know whether Theorem 1.8 is true for the sequence of lamplighter
graphs on the weighted complete binary trees
\((\Lam(\WB_{\theta,n}))_n\), for all or some \(\theta\geq 3\).
\end{quote}

\begin{figure}[htbp]
\centering
\includegraphics[width=\textwidth]{figures/open_problem_crop.png}
\caption{The source paper's weighted-tree follow-up question, page 5 of the
arXiv PDF.}
\end{figure}

\section*{Definitions}

Let \(\D_k\) denote the \(k\)-th diamond graph, with graph diameter \(2^k\).
For an integer \(\theta\geq 2\), the weighted complete binary tree
\(\WB_{\theta,n}\) is the depth-\(n\) complete binary tree in which every edge
whose lower endpoint has level \(j\) has weight \(\theta^{n-j}\). Thus the two
top edges have weight \(\theta^{n-1}\), and each principal subtree below a top
edge is isometric to \(\WB_{\theta,n-1}\).

For a finite weighted tree \(T\), the lamplighter metric is
\[
 d_{\Lam(T)}((A,x),(B,y))
 =
 \operatorname{TSP}_T(x,A\symdiff B,y)+|A\symdiff B|,
\]
where \(A,B\subseteq V(T)\), \(x,y\in V(T)\), and
\(\operatorname{TSP}_T(x,S,y)\) is the shortest weighted walk starting at \(x\),
ending at \(y\), and visiting every vertex of \(S\).

\section*{Proof Intuition}

The subdivided trees \(W_{\theta,n}\) have many vertices along a long top-scale
edge, so the source paper needs a delicate localization argument inside a long
path. The weighted trees have no such interior vertices. This makes the top
edge a rigid cut.

If two lamplighter states are closer than the weight of a tree edge \(e\), then
they cannot be on different sides of \(e\), and they cannot differ by any lamp
on the other side of \(e\). Hence a connected graph whose edges all map below
that scale is frozen with respect to the cut \(e\). Applied to the two top
edges of \(\WB_{\theta,n}\), a minimal-depth diamond embedding must have some
diamond edge whose image has length at least \(\theta^{n-1}\).

But when \(\theta\geq 3\), the whole lamplighter graph has diameter at most a
constant times \(\theta^n\). The image of the two endpoints of \(\D_k\) then
forces \(2^k\lesssim D\theta\), where \(D\) is the distortion. Thus \(k\) is
bounded for each fixed \(D\), which rules out equi-bi-Lipschitz embeddings of
the entire diamond sequence.

\section*{Main Theorem}

\begin{theorem}
\label{thm:weighted-no-diamonds}
For every integer \(\theta\geq 3\), the sequence of diamond graphs
\((\D_k)_{k\geq 1}\) does not equi-bi-Lipschitzly embed into the sequence
\((\Lam(\WB_{\theta,n}))_{n\geq 1}\).
\end{theorem}

\begin{lemma}[Diameter bound]
\label{lem:diameter}
For every \(\theta\geq 3\) and \(n\geq 1\),
\[
  \diam \Lam(\WB_{\theta,n}) \leq 6\theta^n .
\]
\end{lemma}

\begin{proof}
Let \(T=\WB_{\theta,n}\). In a tree, a walk visiting any prescribed lamp
difference set never needs more than twice the total edge weight of \(T\).
Therefore
\[
  \diam\Lam(T)
  \leq 2\sum_{j=1}^n 2^j\theta^{n-j} + (2^{n+1}-1).
\]
Since \(\theta\geq 3\),
\[
  \sum_{j=1}^n 2^j\theta^{n-j}
  \leq \theta^n\sum_{j=1}^{\infty}\left(\frac2\theta\right)^j
  =
  \frac{2}{\theta-2}\theta^n
  \leq 2\theta^n,
\]
and \(2^{n+1}-1\leq 2\theta^n\). Hence
\[
  \diam\Lam(T)\leq 4\theta^n+2\theta^n=6\theta^n .
\]
\end{proof}

\begin{lemma}[Top-edge freezing]
\label{lem:freezing}
Let \(T=\WB_{\theta,n}\), \(n\geq 1\), and let \(e\) be one of the two top
edges. Removing \(e\) separates \(T\) into the principal subtree \(T_e\) below
\(e\) and its complement. If two lamplighter states
\((A,x),(B,y)\in\Lam(T)\) satisfy
\[
  d_{\Lam(T)}((A,x),(B,y)) < \theta^{n-1},
\]
then \(x\) and \(y\) lie on the same side of \(e\). Moreover, if they lie
outside \(T_e\), then \(A\cap T_e=B\cap T_e\), while if they lie inside
\(T_e\), then \(A\setminus T_e=B\setminus T_e\).
\end{lemma}

\begin{proof}
The edge \(e\) has weight \(\theta^{n-1}\). Any weighted walk from one side of
\(e\) to the other must traverse \(e\), so it has length at least
\(\theta^{n-1}\). Thus the lamplighter positions cannot be on different sides.
If the positions are on one side but \(A\symdiff B\) contains a vertex on the
other side, any walk that visits the changed lamp must also traverse \(e\),
again giving length at least \(\theta^{n-1}\). This proves the claim.
\end{proof}

\begin{lemma}[Minimal-depth lower scale]
\label{lem:minimal-depth}
Fix \(D\geq 1\), \(k\geq 1\), and \(\theta\geq 3\). Suppose \(n\) is minimal
among all depths for which there is a \(D\)-bi-Lipschitz embedding
\(\phi:\D_k\to \Lam(\WB_{\theta,n})\). If \(\lambda>0\) is a scaling factor for
\(\phi\), so that
\[
  \lambda d_{\D}(u,v)
  \leq d_{\Lam(\WB_{\theta,n})}(\phi(u),\phi(v))
  \leq D\lambda d_{\D}(u,v),
\]
then
\[
  D\lambda \geq \theta^{n-1}.
\]
\end{lemma}

\begin{proof}
Assume instead that \(D\lambda<\theta^{n-1}\). Since every edge of \(\D_k\)
has length \(1\), adjacent vertices of \(\D_k\) map to lamplighter states at
distance \(<\theta^{n-1}\).

Apply Lemma~\ref{lem:freezing} to each of the two top edges of
\(\WB_{\theta,n}\). Because \(\D_k\) is connected, for each top edge the side
containing the lamplighter position is constant throughout \(\phi(\D_k)\).

If all lamplighter positions lie in one principal subtree \(T_e\), then
Lemma~\ref{lem:freezing} also says that the lamp configuration outside \(T_e\)
is constant over the image. Toggling this fixed outside set is an isometry of
the lamplighter graph, since symmetric differences are unchanged. After this
isometry the image lies in \(\Lam(T_e)\). But \(T_e\) is isometric to
\(\WB_{\theta,n-1}\), contradicting the minimality of \(n\).

The only remaining possibility is that all lamplighter positions lie outside
both principal subtrees, hence at the root. Then the lamp configuration in each
principal subtree is constant over the image. After toggling those two fixed
lamp sets, only the root lamp can vary, so the image has at most two points.
This is impossible for an injective bi-Lipschitz image of \(\D_k\) with
\(k\geq 1\).

Thus \(D\lambda\geq\theta^{n-1}\).
\end{proof}

\begin{proof}[Proof of Theorem~\ref{thm:weighted-no-diamonds}]
Suppose, toward a contradiction, that the sequence \((\D_k)\) embeds
equi-bi-Lipschitzly into \((\Lam(\WB_{\theta,n}))\) with distortion at most
\(D\). For each \(k\), choose \(n\) minimal among all target depths admitting
such an embedding, and let \(\lambda\) be its scaling factor.

Let \(s,t\) be the two original endpoints of \(\D_k\), so
\(d_{\D}(s,t)=2^k\). By Lemma~\ref{lem:diameter},
\[
  \lambda 2^k
  \leq d_{\Lam(\WB_{\theta,n})}(\phi(s),\phi(t))
  \leq 6\theta^n .
\]
By Lemma~\ref{lem:minimal-depth}, \(\theta^{n-1}\leq D\lambda\). Combining the
two estimates gives
\[
  2^k
  \leq
  \frac{6D\theta^n}{\theta^{n-1}}
  =
  6D\theta .
\]
For fixed \(D\) and \(\theta\), this bounds \(k\), contradicting the existence
of uniformly distorted embeddings for all diamond levels.
\end{proof}

\section*{Verification Notes}

The accompanying script \texttt{code/verify\_scale\_bounds.py} checks the elementary
finite-sum bound
\[
2\sum_{j=1}^n2^j\theta^{n-j}+(2^{n+1}-1)\leq 6\theta^n
\]
for a finite range of \(\theta,n\), and checks the resulting threshold
\(2^k>6D\theta\). The proof above gives the general symbolic argument.

\section*{Limitations and Novelty Check}

This packet answers only the second follow-up question on page 5 of the source
paper: the no-diamonds theorem for weighted complete binary trees. It does not
answer the separate question whether \((\Lam(W_{\theta,n}))_n\) for
\(\theta\geq 3\) embeds into every nonsuperreflexive Banach space.

The run indexes were searched for arXiv:2509.25421, lamplighter graphs,
weighted binary trees, diamonds, and the adjacent BMSZ Problem 7.1 family. The
only relevant prior attempt was an older stretched-tree push for arXiv:1902.07098;
that target is answered by the present source paper and is not duplicated here.
A bounded web search on July 18, 2026 for exact phrases involving
``Lam(WB)'', ``weighted complete binary trees'', lamplighter graphs, and
diamonds found the source paper but no later answer to the weighted-tree
follow-up.

\section*{References}

\begin{itemize}
\item C. Melby and B. Randrianantoanina, \emph{Bi-Lipschitz embedding
properties of lamplighter graphs on weighted and unweighted trees},
arXiv:2509.25421.
\item F. Baudier, P. Motakis, T. Schlumprecht, and A. Zsak,
\emph{On the bi-Lipschitz geometry of lamplighter graphs}, Discrete Comput.
Geom. 66 (2021), 203--235.
\item W. B. Johnson and G. Schechtman, \emph{Diamond graphs and
super-reflexivity}, J. Topol. Anal. 1 (2009), 177--189.
\end{itemize}

\section*{Human Review Recommendation}

Check the cut-freezing Lemma~\ref{lem:freezing} against the source paper's
normalization of weighted lamplighter metrics. Under the stated TSP-plus-lamp
metric, the proof is insensitive to all constants except the strict comparison
\(D\lambda<\theta^{n-1}\), and equality is handled by the contrapositive in
Lemma~\ref{lem:minimal-depth}.

\end{document}
